\paragraph{What the reported exponents can show.} The exponents 1.49, 1.65 and 3.43 of \cite{stringer2019} are descriptive statistics of particular stimulus sets over particular rank windows. They are consistent with the bound, since codes that satisfy it can produce them. They cannot show that the code satisfies it, nor the stronger conclusions that it is ``as high-dimensional as possible'' or that the asymptotic exponent is close to $1 + 2/d$ (the fitted exponents are ``near but above the lower bounds''; the objection is to reading them as the asymptotic exponent), because codes that violate the bound, or sit exactly at it, produce the same window exponents. Border codes with 8,704 neurons reach 1.49 in \FinReachEight{} of the six 8D sets in the unwhitened coordinates (in whitened ones none does with infinitely many neurons) and 1.65 in \FinReachAllFour{} four 4D sets; codes below the border exceed the bound in some settings (Section~\ref{sec:results}); and on the torus a differentiable and a non-differentiable code can have finite-sample spectra that differ by less than any given amount at every rank, for every stimulus set (Corollary~\ref{cor:torus}). Conversely, a window exponent below $1 + 2/d$ would not show that a code is non-differentiable, since smooth codes with short length scales give such values. The ordinal prediction, larger exponents for lower $d$, is reproduced by codes of fixed smoothness, so its confirmation does not bear on smoothness either. A window exponent measures how the code's variance is spread over the modes the stimulus set resolves, at scales set by the tuning. This bears on the reanalysis of Davidovich and Roudi \cite{dr2022}, who found the 8D exponents above the bound ``with a reasonable margin'' and wrote that the uncertainty of the fits ``does not affect the conclusions regarding the smoothness of the population response to low-dimensional stimuli'': their result concerns the fitted exponent, and our computations say that even an exactly known window exponent does not decide smoothness. The broken-power-law tail exponent is likewise set by the assumed continuation beyond the resolved ranks: a tail that departs from the assumed form changes the trace, the fit absorbs this into $\alpha_2$, and the misfit has little power to detect it, with nominal flag rates from 0 to \MemeFlagAnyMax{} depending on the reference distribution and the covariance estimate. Pospisil and Pillow show that MEME recovers eigenvalues beyond the rank of the data when the parametric form is right (their Supplementary Fig.~S1, ``Demonstration that MEME estimator can infer eigenvalues beyond the rank of data''); they also note that the form of the tail ``remains unclear''.

\paragraph{What a test would need.} One route is many more stimuli per dimension, with the window moved to higher ranks as well: as $P$ grows, the fixed ranks 11--500 converge to the operator eigenvalues at those ranks, which for long tuning length scales remain pre-asymptotic, so more stimuli alone do not bring the window into the asymptotic regime. How many would be needed depends on the unknown length scale. A heuristic shows the scale of the problem: the nearest-neighbour spacing of $P$ stimuli in $d$ dimensions shrinks as $P^{-1/d}$, so resolving scales ten times shorter than 2,800 stimuli resolve now takes of the order of $10^d$ times as many stimuli. At $d = 1$ the obstacle is the number of directions: with 32 the window ends at the sampling frequency. The other route is a parametric model of the code, for example a kernel family such as \eqref{eq:matern} fitted to how response similarity falls with stimulus distance, with the bound read off the fitted smoothness. Such a test holds within the model, depends on the metric in which the code is assumed stationary, and meets the same resolution limit, because smoothness is decided at the shortest distances the stimulus set contains. Either route needs its assumptions stated in advance, and the geometry of each stimulus set, not only the nominal $d$, must enter the analysis.

\paragraph{Limitations.} The codes placed on the stimulus coordinates are isotropic, stationary Gaussian-process codes in the Euclidean metric of the image principal components, or of their whitened version. Real codes are neither isotropic nor stationary, and the metric in which the neural code is smooth is unknown; the family shows what the window exponent can do at fixed smoothness, not how visual cortex behaves. Proposition~\ref{prop:window} is proved for bounded eigenfunctions, and for the Mat\'ern codes on $\mathbb{R}^d$ we rely on computation. The identifiability computations are noise-free; estimation error comes on top. The MEME computations use one implementation with stated deviations from the authors' code, one noise model with fixed neurons and illustrative spectra, and for the Mat\'ern codes weights and a misfit threshold calibrated on a different spectrum; the size of the shifts in $\alpha_2$ depends on these choices, their existence does not. The stimulus-count dependence rests on 4 subsets per size. We make no statement about the exponents of the 8D, 4D or grating recordings. The bound of \cite{stringer2019} is a theorem, and nothing here questions it; the question is only whether data of this size can show that a code satisfies it.
